\section{Composition setup}
\label{sec:composition_setup}

The deployment-safety theorem of \S\ref{sec:main_theorem} composes results from
six papers in the GFM sequence.  This section sets up the composition: what each
source paper contributes, the intensive-vs-extensive distinction that makes the
composition tractable, the notation reconciliation across source papers, and
the composition challenges this paper must address.

\subsection{Source-paper contributions as building blocks}
\label{sec:source_blocks}

The composition uses six load-bearing source-paper results.  Each is
intensive: stated in per-capability, per-channel, or per-event terms
rather than in terms of absolute capability magnitude.  This intensity
is what makes the composition's capability-magnitude-independent
bound possible.

\textbf{The Horizon-Aware paper~\citep{lasser2026horizon}: The anti-monopolar property.}
\citep[Proposition~6]{lasser2026horizon} establishes that under a
discounted objective with positive cross-substrate cooperative novelty
$\rext > 0$, full capability domination is anti-maximizing:
\[
  \Vdisc^{\mathrm{div}} - \Vdisc^{D} \;=\;
  \frac{\rext}{1 - \gamma} - \Delta_0,
\]
where $\Vdisc^{\mathrm{div}}$ and $\Vdisc^{D}$ are the discounted
values of the diversity-maintaining and domination strategies, and
$\Delta_0$ is the immediate $\volL$-change from domination.
Diversity strictly dominates for $\gamma > \gamma^* = 1 -
\rext/\Delta_0$ (when $\Delta_0 > 0$) or for all $\gamma \in (0,1)$
(when $\Delta_0 \leq 0$). The corollary on strategy-dependent
internal rate identifies the breakdown regime $\Delta r_K \geq
\rext$, which this paper's safe-region analysis must constrain.

\textbf{The Exogenous-Verification paper~\citep{lasser2026exo}: Verification infrastructure.}
\citep{lasser2026exo} supplies the operational substrate for ledger-observable
invariants: substrate-exclusive algorithmic witnesses (Definition of
Algorithmic Witness), Pedersen commitments (computationally binding, perfectly
hiding under discrete-log assumption), an append-only verification ledger
replicated across substrate-exclusive nodes, and a governance fork protocol
requiring cross-substrate supermajority for evaluation-protocol acceptance. The
behavioral consistency monitor (\citep[Proposition~6]{lasser2026exo}) provides
SPRT detection with expected detection time $\mathbb{E}[T_{\mathrm{detect}}]
\leq A/\delta$ where $A$ is the SPRT upper threshold (Wald's sequential
test~\citep{wald1947sequential}) and $\delta$ is the KL
divergence~\citep{kullback1951information} between baseline and alternative
behavioral distributions.

\textbf{The Phase-Redundancy paper~\citep{lasser2026phase}: Phase boundary dynamics.}
\citep{lasser2026phase} provides a Lyapunov function $\Lyap(\Wm_t) = \sum_k w_k
\epsilon_k(t)^2$ on world-model error and a phase boundary theorem
characterizing when the coupled $(P, \Wm)$ system stays in the self-correcting
basin versus enters the monopolar absorbing state. The critical surface
(\citep[Equation 7 (critical surface)]{lasser2026phase}) is expressed in
intensive quantities:
\[
  \rS \cdot \rhomincross + \alpha_{\min} c_V \underline{\nu}
  \;>\; \rW \cdot \rsub + \dmax \cdot
  \frac{\Delta\rho_{\mathrm{avg}}}{\Lyap_{\mathrm{avg}}}
\]
where $\rhomincross$ is the cross-substrate redundancy minimum, $\rsub$ is the
subsumption frequency, and $\rS, \rW, \dmax$ are intensive rate constants. The
metastable lifetime $\tau_{\mathrm{meta}} \gtrsim C/I_k$ under partial
endogenous correction (B1$'$) provides a lower bound on the cascade time that
this paper's lead-time guarantee uses.

\textbf{The Revealed-Sacrifice paper~\citep{lasser2026paper8a}: B-to-C lower bound.}
\citep{lasser2026paper8a} establishes the aggregate B-to-C lower bound theorem:
$\volRwin(Y_n) \geq \Delta\volL(X_n)$ per revealed-sacrifice event, with
monotone accumulation ($\volRlower(E') \geq \volRlower(E)$ for $E \subseteq
E'$). The B-to-C ratio $\betaL = \volRlower / \volL \in [0, 1]$ is the
ledger-derived measure of how much realized exercise has been witnessed against
possessed capability volume.

\textbf{The Need-Sufficiency paper~\citep{lasser2026paper8b}: HHI and gap decomposition.}
\citep{lasser2026paper8b} establishes the trade-flow Herfindahl-Hirschman Index
$\HHI$ as a wireheading-consistent concentration signal, Schur-convex in the
majorization order, with the Part-A category-flow variant third-party
observable from public committed events plus public S1-admissibility labelling.
The gap decomposition partitions the B-to-C gap into five cells (restricted,
covered, dormant, residual, boundary-residual) with classification computable
in near-linear $\widetilde{O}(|P|+E+M)$ time.

\textbf{The Microfoundation paper~\citep{lasser2026micro}: Static Goodhart bound.}
\citep{lasser2026micro} establishes the Lipschitz transfer Goodhart
bound: for $\Pproxy = \volL$ as proxy and $\Ttruth = \volRwin$ as
operational truth (under stance S0), and Lipschitz alignment
property $g$:
\[
  |g(\Ttruth) - g(\Pproxy)| \;\leq\; \mathrm{Lip}(g) \cdot \epsnonres,
\]
where $\epsnonres$ is the non-residual proxy-truth gap on the operationally
active subspace $P^{\mathrm{act}}$. The four-channel decomposition (observation
density, attestation quality, individuation discipline, bundle decomposition)
governs the gap's evolution. \citep[Microfoundation]{lasser2026micro}'s
universal model-free Concentration-Gap~Conjecture (optimization pressure
correlates with gap exploitation) remains open as a research-program
target. The companion paper~\citep{lasser2026foundations} discharges a
\emph{scoped} version under four structural conditions
(Assumption~\ref{ass:concgap_conditions}: REP, DISP, COAL, EMB);
this paper's $I_5$ (HHI ceiling) provides the operational HHI
threshold under which the scoped discharge binds.

\subsection{The intensive-vs-extensive distinction}
\label{sec:intensive_extensive}

A bound is \emph{intensive} if it is stated in terms of per-element quantities
(per-capability, per-channel, per-event) whose magnitude does not grow with the
size of the underlying capability poset $P$ or agent population. A bound is
\emph{extensive} if it scales with $|P|$ or with capability count.

The capability-magnitude-independent property of this paper's main theorem
requires the bound on Goodhart slack to be intensive in $|P|$. If any
source-paper bound were extensive --- for example, if $\epsnonres$ grew as a
sum over capabilities rather than as a per-capability supremum --- the
composed bound would inherit that extensive scaling, and the deployment claim
would weaken with capability magnitude.

Each source-paper result above is intensive by construction:
\begin{itemize}
\item \citep[Horizon~Aware]{lasser2026horizon}'s $\rext$ and $\Delta r_K$ are growth rates, not
  capability counts.
\item \citep[Exogenous~Verification]{lasser2026exo}'s SPRT detection rate $A/\delta$ depends on KL
  divergence and threshold, not on capability count.
\item \citep[Phase~Redundancy]{lasser2026phase}'s critical surface uses minima and averages
  ($\rhomincross$, $\Delta\rho_{\mathrm{avg}}$), not sums.
\item \citep[Revealed~Sacrifice]{lasser2026paper8a}'s $\betaL$ is a ratio in $[0,1]$.
\item \citep[Need~Sufficiency]{lasser2026paper8b}'s $\HHI$ is a concentration index in $[1/n, 1]$ where
  $n$ is the number of categories.
\item \citep[Microfoundation]{lasser2026micro}'s $\epsnonres$ is a sup-norm on $P^{\mathrm{act}}$, not
  an average over capability count.
\end{itemize}

The composition challenge is to show that the \emph{combination} of
these intensive bounds remains intensive. Section~\ref{sec:lemmas}
addresses this via Lemma~\ref{lem:1} (intensive composition under
co-evolution); the proof handles \citep[Microfoundation]{lasser2026micro}'s
Composition~Proposition~1 positive-part error terms explicitly.

\subsection{Notation reconciliation}
\label{sec:notation}

The source papers use different conventions for the same underlying
quantities; this paper standardizes on a single set of symbols
across the composition.  The full notation reconciliation
(source-paper symbol $\to$ This paper's symbol $\to$ meaning) is
collected in Appendix~\ref{app:notation}, Table~\ref{tab:notation}.

The most consequential reconciliations are:
\begin{itemize}
\item \citep[Phase Redundancy]{lasser2026phase}'s $\meff$ (nominal
  substrate count) is upgraded to this paper's $\mindep$
  (failure-correlation-independent substrate count).  Nominally
  distinct substrates may share failure modes (skeleton substrates);
  This paper's safe-region calculation requires genuine independence.
\item \citep[Microfoundation]{lasser2026micro}'s $P$ and $T$ are
  written as $\Pproxy = \volL$ and $\Ttruth = \volRwin$ throughout
  this paper to prevent confusion with $P$ as a poset and $T$ as a
  time variable.
\item \citep[Horizon Aware]{lasser2026horizon}'s discounted value
  $V^\gamma$ is written as $\Vdisc$ to make the discount factor an
  explicit subscript rather than a superscript that conflicts with
  strategy labels ($\Vdisc^{\mathrm{div}}$, $\Vdisc^D$).
\end{itemize}

\subsection{Composition challenges}
\label{sec:composition_challenges}

The source-paper results do not compose mechanically. Three challenges drive
the technical content of \S\ref{sec:lemmas} and \S\ref{sec:main_theorem}.

\textbf{Challenge 1: Co-evolution of channels.}  \citep[Microfoundation]{lasser2026micro}'s
Composition~Proposition~1 establishes that the four-channel
decomposition composes \emph{exactly} in the sequential-intervention
regime but only \emph{approximately} under co-evolution, with
positive-part error terms. This paper's intensive composition lemma
(Lemma~\ref{lem:1}) propagates these error terms through the
deployment-safety bound. The error terms are bounded but non-zero;
the deployment claim must accommodate them.

\textbf{Challenge 2: Lyapunov-to-Goodhart bridge.}  \citep[Phase~Redundancy]{lasser2026phase}'s
Lyapunov function tracks world-model error, and \citep[Microfoundation]{lasser2026micro}'s $\epsnonres$
tracks proxy-truth gap on the active subspace. These quantities are
distinct: $\Lyap$ is per-dimension squared error in $\Wm$, while
$\epsnonres$ is sup-norm gap in the proxy-truth comparison.
Lemma~\ref{lem:2} establishes the quantitative bridge:
$\Lyap < \epsilon_{\mathrm{safe}}$ implies $\epsnonres < f(\epsilon_{\mathrm{safe}})$
with $f$ explicit and intensive in capability magnitude.

\textbf{Challenge 3: HHI-to-pressure operationalization.}  \citep[Microfoundation]{lasser2026micro}'s
universal model-free Concentration-Gap~Conjecture (optimization
pressure correlates with gap exploitation) is unproved. The
companion paper~\citep{lasser2026foundations} closes the
deployment-relevant gap by proving a \emph{scoped} version
under four structural conditions
(Assumption~\ref{ass:concgap_conditions}: REP, DISP, COAL,
EMB), and Layer~1 of Theorem~\ref{thm:deployment_safety} now
binds through that scoped structural discharge rather than
through direct conditioning on the unproved universal conjecture.
This paper's $I_5$ (HHI ceiling) supplies the HHI threshold
side of the discharge: under (REP)+(DISP)+(COAL)+(EMB) plus
$\HHI < H^*$, \citep[Theorem~2]{lasser2026foundations} delivers
the proxy-truth Goodhart-slack bound. Lemma~\ref{lem:3_forward}
formalizes the forward direction
($\HHI > H^* \Rightarrow$ trade-flow concentration prerequisites
of the optimization-pressure regime); the reverse direction
required by Layer~1 is supplied by the scoped structural
theorem. The universal conjecture remains a research-program
target deferred
indefinitely (\S\ref{sec:disc_indefinite}); the deployment
claim does not depend on it directly.

\subsection{What this paper supplies on top of source-paper machinery}
\label{sec:paper10_new}

This paper's substantive new content (beyond composition of
inherited results):
\begin{enumerate}
\item \textbf{The eleven operational invariants}
  ($I_1$--$I_{11}$, \S\ref{sec:invariants}), with definitions,
  threshold semantics, and ledger-observable measurement procedures.
\item \textbf{The five-part Lemma~5 family} (\S\ref{sec:lemmas}):
  substrate floor, channel-restricted detection, minimax static
  tightening, lead-time tail bound, environment-side witness
  extension.
\item \textbf{The three-layer deployment-safety theorem}
  (\S\ref{sec:main_theorem}): static safe region, detection-and-
  correction, named residuals.
\item \textbf{The canonical tripartite substrate identification}
  (\S\ref{sec:substrate_anchoring}): Human + AI +
  Formal-Operational, with the cooperative-anchoring property.
\item \textbf{Three additional invariants for cooperative-anchoring
  evasions} ($I_9$, $I_{10}$, $I_{11}$): substrate-exclusivity
  observability, coverage/materiality, latency bounds.
\item \textbf{The deployment-tooling specification}
  (\S\ref{sec:deployment_tooling}): what an operator must
  instrument to invoke the theorem.
\end{enumerate}

The composition is non-trivial: the five-part Lemma~5 family
carries internal asymmetries (Lemma~5a's substrate floor relies on
pairwise-additivity assumptions distinct from the failure-
correlation independence used elsewhere; Lemma~5c's safe region is
binding only in the cooperative-overlap regime), and the
cooperative-anchoring property of \S\ref{sec:substrate_anchoring}
holds in a deliberately narrow form rather than as a general
inner-alignment guarantee.  The result is a bounded but defensible
structural claim, not a sweeping safety guarantee.
\S\ref{sec:scope} of the introduction stated the explicit residuals;
\S\ref{sec:discussion} returns to them after the formal apparatus
is in place.
